% \iffalse meta-comment
%
% hit-hyperlink.dtx 是各个类共用的 hyperref 配置
%
% \fi
%
% \section{共享超链接配置}
%
% \changes{v3.2a}{0000/00/00}{hyperref 的配置抽出来给各个类共用，原先每个类各写
%   一遍，十三行逐字节相同}
% PDF 书签、链接边框、url 字体这些跟学校规范无关，是模板整体的取舍，三个类原先
% 各写一遍且一字不差。
%
% \cs{hit@setup@hyperlink} 只做配置，由各个类在原来那段代码所在的位置调用，
% 所以生成物里的加载顺序不变。\pkg{natbib} 的引用样式补丁没有一起搬——那是
% \file{hithesisbook.cls} 独有的，也不属于超链接。
%
%    \begin{macrocode}
%<*shared-hyperlink>
%    \end{macrocode}
%
% \begin{macro}{\hit@setup@hyperlink}
% \texttt{colorlinks=false} 加 \texttt{pdfborder=0 0 0}：屏幕上不上色也不画框，
% 打印出来与正文一致。\texttt{CJKbookmarks} 让中文书签不乱码，
% \texttt{bookmarksopenlevel=1} 只展开到章。
%    \begin{macrocode}
\cs_new_protected:Npn \hit@setup@hyperlink
  {
\RequirePackage{hyperref}
\hypersetup{%
  CJKbookmarks=true,
  linktoc=all,
  bookmarksnumbered=true,
  bookmarksopen=true,
  bookmarksopenlevel=1,
  breaklinks=true,
  colorlinks=false,
  plainpages=false,
  pdfborder=0~0~0
}
\urlstyle{same}
  }
%    \end{macrocode}
% \end{macro}
%
%    \begin{macrocode}
%</shared-hyperlink>
%    \end{macrocode}
%
% \Finale
